% ============================================
% LUC-06
% Gerado de molambudos.md — build_miolo.py
% ============================================
\fragtitle{luc}{L}{LUC}{06}{O Desaparecimento}
\markboth{LUC-06 --- O Desaparecimento}{}

\noindent \textit{O ciclo não tem reversa}

\noindent \textbf{Maio de 2026 --- Belo Horizonte}

\vspace{0.5em}\noindent\rule{\textwidth}{0.5pt}\vspace{0.5em}

\noindent O último e-mail de Dra. Lúcia Mendes foi enviado em 12 de maio de 2026, às 2:47 da manhã.

\noindent O destinatário era o Dr. Ronaldo Torres, do Arquivo Público Mineiro.

\vspace{0.5em}\noindent\rule{\textwidth}{0.5pt}\vspace{0.5em}

\noindent \textit{\textquotedblleft{}Ronaldo,}

\noindent \textit{Preciso que você guarde este relato. Se algo acontecer comigo, quero que fique registrado.}

\noindent \textit{Li o diário completo do paciente 1.260. São 62 anos de registro --- de 1917 a 1979. As últimas páginas foram escritas depois da morte dele. Eu sei que parece impossível. Mas as páginas estão lá. A tinta está lá. A caligrafia é a mesma.}

\noindent \textit{Joaquim Antônio Correia não estava louco.}

\noindent \textit{Ele estava possuído.}

\noindent \textit{Eu sei que isso soa absurdo vindo de uma psiquiatra forense. Eu passei 15 anos diagnosticando traumas, dissociações, psicoses. Eu acreditava que tudo tinha uma explicação neurológica ou psicológica.}

\noindent \textit{Agora acredito que existem coisas que a psicologia não explica.}

\noindent \textit{Li o diário. Três dias depois, comecei a sonhar com a vala. Uma semana depois, senti o cheiro. Duas semanas depois, meu olho direito começou a clarear.}

\noindent \textit{O mesmo olho. O mesmo amarelo.}

\noindent \textit{Fiz exames. Tomografia. Ressonância. Sangue. Não há explicação médica.}

\noindent \textit{Fui ao Colônia. Encontrei um caderno de Joaquim escondido numa cela --- dos anos de mutismo. Encontrei um velho que trabalhou lá nos anos 50. Ele me disse: 'A senhora leu o diário. Agora é sua vez.'}

\noindent \textit{Perguntei a vez do quê. Ele disse: 'De hospedar.'}

\noindent \textit{Estou com medo, Ronaldo. O medo não passa pela razão. Fica no corpo. Um medo que vem do estômago. Um medo que tem cheiro doce.}

\noindent \textit{Hoje, olhei no espelho e vi o olho amarelo me olhando de volta.}

\noindent \textit{Não era o meu olho.}

\noindent \textit{Era o olho dele.}

\noindent \textit{Olho de Joaquim.}

\noindent \textit{Ou daquilo que estava dentro dele.}

\noindent \textit{Vou entregar o diário no arquivo amanhã. Vou colocar num pano preto e deixar na recepção. Depois disso, vou sumir por um tempo.}

\noindent \textit{Não procure por mim.}

\noindent \textit{Se alguém perguntar, diga que o diário foi devolvido.}

\noindent \textit{E que o ciclo continua.}

\noindent \textit{Lúcia\textquotedblright{}}

\vspace{0.5em}\noindent\rule{\textwidth}{0.5pt}\vspace{0.5em}

\noindent O diário foi encontrado na recepção do Arquivo Público Mineiro em 13 de maio de 2026, envolto num pano preto.

\noindent Dra. Lúcia Mendes não foi vista em público desde então. Os registros de movimentação do arquivo, porém, mostram que ela voltou ao prédio em julho, sempre à noite, e que assinou a retirada do objeto em 28 de julho.

\noindent Seu apartamento foi encontrado vazio. A cama estava arrumada. O café estava frio na xícara. O computador estava aberto na página de pesquisa: \textit{\textquotedblleft{}Como passar o diário adiante sem ler?\textquotedblright{}}

\noindent Não havia resposta.

\noindent Não há resposta.

\vspace{0.5em}\noindent\rule{\textwidth}{0.5pt}\vspace{0.5em}

\noindent \textit{--- Eu deveria ter queimado o diário. Mas não queimei.}
\textit{--- Eu deveria ter parado de ler. Mas não parei.}
\textit{--- Eu deveria ter fugido antes do olho amarelo. Mas já era tarde quando senti o cheiro.}

\noindent Se você está lendo isto, você já sentiu o cheiro também.
Você já sonhou com a vala.
Você já viu o número 1.260 em tudo.

\noindent \textbf{PARE AGORA.}

\noindent Feche o livro. Devolva. Não termine este capítulo.

\noindent Eu estou escrevendo isto de um computador de lan house em uma cidade que não é a minha. Meu olho direito já não enxerga mais como antes. Tudo tem uma tonalidade amarelada. Como se eu estivesse vendo o mundo através de um olho que não é meu.

\noindent Se você continuar, vai acontecer com você também.

\noindent A criatura não está no diário.
A criatura está na \textit{leitura}.
Ela existe no espaço entre a página e o cérebro.

\noindent Se você está lendo, você já abriu a porta.
Mas ainda pode fechar.

\noindent \textit{--- Lúcia Mendes, paciente 1.262. Último aviso.}

\vspace{0.5em}
\noindent\textit{\textbf{↪ Rotas:} \rota{LUC-07} \rotasep \rota{CONT-05}}